\documentclass[10pt,a4paper]{amsart}
\usepackage[margin=19mm]{geometry}
\usepackage{amsmath,amssymb,mathtools}
\usepackage[scr=rsfs,cal=euler]{mathalfa}
\usepackage{xfrac}
\usepackage[unicode,hidelinks,bookmarks=false]{hyperref}
\newcommand{\ie}{i.e.\ }
\newcommand{\cf}{cf.\ }
\newcommand{\resp}{respectively\ }
\newcommand{\Subsection}{\subsection}
\providecommand{\nobreakdash}{\nobreak\hbox{-}}
\setlength{\emergencystretch}{2em}
\multlinegap0pt
\usepackage{bm}
\usepackage{bbm}
\usepackage{dsfont}
\usepackage{xspace}
\usepackage{cancel}
\usepackage{mathtools}
\usepackage[shortlabels]{enumitem}
\usepackage{amsthm}
\makeatletter
\@ifundefined{theorem}{\newtheorem{theorem}{Theorem}}{}
\@ifundefined{proposition}{\newtheorem{proposition}[theorem]{Proposition}}{}
\makeatother
\makeatletter
\newcommand{\CC}{\mathcal{C}}
\newcommand{\MM}{\mathcal{M}}
\newcommand{\bk}[1]{\mathbf{#1}}
\newcommand{\ind}{\mathrm{Ind}}
\newcommand{\kk}{\mathbb{K}}
\expandafter\ifx\csname proof\endcsname\relax
\renewenvironment{proof}{{\bf \noindent Proof.}}{\qed}
\fi
\newcommand{\wt}{\widetilde}
\makeatother
\title[Graded Betti numbers of generalized split--join graphs and a]{Graded Betti numbers of generalized split--join graphs and applications}
\author{Bilal Ahmad Rather}
\date{Source: arXiv:2606.24822v1 (CC BY 4.0)}
\begin{document}
\maketitle

\noindent\emph{Source and licence.} Graded Betti numbers of generalized split--join graphs and applications. Version arXiv:2606.24822v1 (CC BY 4.0), distributed under the Creative Commons Attribution 4.0 International licence, as recorded in the arXiv licence metadata for this exact version and shown on the abstract page https://arxiv.org/abs/2606.24822v1. Full rights notice: licenses/arxiv-2606.24822-CC-BY-4.0.txt.\par\smallskip

\noindent\textit{Editorial scope.} Counted unit: the Theorem whose source label is \texttt{thm:split-linear} in the article (source0.tex lines 574--587, source label \texttt{thm:split-linear}), delivered together with the article's own complete original proof of it (source0.tex lines 589--643, one \texttt{proof} environment of 55 lines). No mathematics is written, repaired or re-derived: statement and proof are the article's own text line for line. Required context retained uncounted: prop:induced-models (lines 280--309). Every definition, display and label of those lines is retained unchanged. Omitted from this package and not redistributed: 1--279, 310--573, 588--588, 644--1687. Nothing inside a retained excerpt is omitted or altered: every retained body is delivered exactly as the article's lines are written, so no omission receipt of any kind accompanies this package. Citations: the retained excerpts carry no citation command at all, so no bibliography entry of the article is transcribed and no reference is redistributed. Reference adaptation: none was needed and none was made; every reference inside the delivered unit resolves inside it, so no unresolved reference or citation is produced. Printed slips kept literal: no printed word, symbol, hypothesis or number of the counted statement or of its proof was altered. Layout and class: the article's own class is \texttt{article} with packages including geometry, amsmath, amssymb, amsthm, amsfonts, mathtools, mathrsfs, bm, booktabs, array, multirow, enumitem, hyperref, tikz; the wrapper is the lane's pinned \texttt{amsart} editorial wrapper at 10pt on A4, which restates the theorem-like environments the retained text needs and adds the article's own macro definitions that the retained text needs (\texttt{\textbackslash CC}, \texttt{\textbackslash MM}, \texttt{\textbackslash bk}, \texttt{\textbackslash ind}, \texttt{\textbackslash kk}, \texttt{\textbackslash wt}), with extra wrapper packages (bm, bbm, dsfont, xspace, cancel, mathtools, enumitem). The delivered numbering is the unit's own. Assets: no retained span contains a figure environment, a graphics-inclusion command, a PDF-page inclusion, a table or a TikZ picture; no image file is copied into this package, the package declares no asset dependency and no image file is redistributed with it. Licence: version arXiv:2606.24822v1 of this article is distributed under the Creative Commons Attribution 4.0 International licence, as recorded in the arXiv licence metadata for this exact version and shown on the abstract page https://arxiv.org/abs/2606.24822v1. A full rights notice is delivered at licenses/arxiv-2606.24822-CC-BY-4.0.txt. Characters: every retained excerpt is pure ASCII, so the delivered engine, XeLaTeX with its default Latin Modern text font, reports no missing character.
	\begin{proposition}\label{prop:induced-models}
		Let $W\subseteq A\sqcup B_1\sqcup\cdots\sqcup B_m$, and let
		$ u=|W\cap A|, $ $ 	w_r=|W\cap B_r|, $ and $t=\#\{r:w_r>0\}. $ 	Then the induced subcomplexes are described as follows.
		
		\begin{enumerate}[label=\textup{(\roman*)},noitemsep]
			\item For $G=\MM(a;\bk{b})$, we have
			$$
			\ind(G)_W \cong
			\begin{cases}
				\langle W\cap A\rangle, & t=0,\\
				D_{w_r}, & u=0,\ t=1,\\
				J(w_{r_1},\dots,w_{r_t}), & u=0,\ t\geq 2,\\
				\langle W\cap A\rangle \sqcup D_{w_r}, & u>0,\ t=1,\\
				\langle W\cap A\rangle \sqcup J(w_{r_1},\dots,w_{r_t}), & u>0,\ t\geq 2.
			\end{cases}
			$$
			
			\item For $G=\CC(a;\bk{b})$, we have
			$$
			\ind(G)_W \cong
			\begin{cases}
				D_u, & t=0,\\
				D_{w_r}, & u=0,\ t=1,\\
				J(w_{r_1},\dots,w_{r_t}), & u=0,\ t\geq 2,\\
				D_u\sqcup D_{w_r}, & u>0,\ t=1,\\
				D_u\sqcup J(w_{r_1},\dots,w_{r_t}), & u>0,\ t\geq 2.
			\end{cases}
			$$
		\end{enumerate}
	\end{proposition}

	\begin{theorem}\label{thm:split-linear}
		For $j\geq 2$,
		$$
		\beta_{j-1,j}(\MM(a;\bk{b}))=
		\sum_{r=1}^m (j-1)\binom{b_r}{j}
		+
		\sum_{u=1}^{\min{a,j-1}}
		\binom{a}{u}
		\sum_{r=1}^m (j-u)\binom{b_r}{j-u}
		+
		\sum_{u=1}^{\min{a,j-2}}
		\binom{a}{u}M_{j-u}(\bk{b}).
		$$
	\end{theorem}

	\begin{proof}
		For $G=\MM(a;\bk{b})$, Hochster's formula implies
		$$
		\beta_{j-1,j}(G)=\sum_{\substack{W\subseteq V(G)\\ |W|=j}}
		\dim_{\kk}\wt H_0(\ind(G)_W;\kk).
		$$
		We therefore classify all $j$-subsets $W$ for which $\ind(G)_W$ is disconnected, and consider the following cases.
		
		\smallskip
		
		\noindent\textbf{Case 1:} If $W\subseteq B_r$ for some fixed $r$, then $\ind(G)_W\cong D_j$, so we have
		$$
		\dim_{\kk}\wt H_0(\ind(G)_W;\kk)=j-1.
		$$
		There are $\binom{b_r}{j}$ such choices inside $B_r$, and hence the total contribution of this case is
		$$
		\sum_{r=1}^m (j-1)\binom{b_r}{j}.
		$$
		
		\smallskip
		
		\noindent\textbf{Case 2:} If $W$ meets $A$ and exactly one block $B_r$, then for $u=|W\cap A|$ with $1\leq u\leq j-1$, so we get $|W\cap B_r|=j-u$. By Proposition~\ref{prop:induced-models}, we have
		$$
		\ind(G)_W\cong \langle W\cap A\rangle \sqcup D_{j-u}.
		$$
		It has $(j-u)+1$ connected components, and hence we have
		$$
		\dim_{\kk}\wt H_0(\ind(G)_W;\kk)=j-u.
		$$
		For fixed $u$, the number of such subsets is $\binom{a}{u}\binom{b_r}{j-u}$, so the total contribution is
		$$
		\sum_{u=1}^{\min{a,j-1}}
		\binom{a}{u}
		\sum_{r=1}^m (j-u)\binom{b_r}{j-u}.
		$$
		
		\smallskip
		
		\noindent\textbf{Case 3:} If $W$ meets $A$ and at least two distinct blocks. Again For $u=|W\cap A|$, with $1\leq u\leq j-2$, and by Proposition~\ref{prop:induced-models}, we have
		$$
		\ind(G)_W\cong \langle W\cap A\rangle \sqcup J(\bk{w}),
		$$
		where $J(\bk{w})$ is connected, since it is the join of at least two nonempty complexes. Hence, $\ind(G)_W$ has exactly two connected components, so we obtain
		$$
		\dim_{\kk}\wt H_0(\ind(G)_W;\kk)=1.
		$$
		For fixed $u$, the number of ways to choose the remaining $j-u$ vertices in at least two blocks is precisely $M_{j-u}(\bk{b})$. SO, in this  case, the total  contribution is
		$$
		\sum_{u=1}^{\min{a,j-2}}
		\binom{a}{u}M_{j-u}(\bk{b}).
		$$
		
		\smallskip
		 No other subset contributes to $\beta_{j-1,j}(\MM(a;\bk{b}))$. Indeed, if $W\subseteq A$, then $\ind(G)_W$ is a simplex, if $W$ is contained in at least two blocks and avoids $A$, then $\ind(G)_W$ is connected, and contributes null. Now, summing the cases, we obtain the stated formula for $\beta_{j-1,j}(\MM(a;\bk{b})).$
	\end{proof}

\end{document}
